Let $G=(V,E)$ be a finite simple undirected connected graph with $|V|=n\ge 1$.
A \emph{path} means a simple path: a sequence of distinct vertices $v_0,v_1,\dots,v_\ell$ with $\ell\ge 1$ and $\{v_{i-1},v_i\}\in E$ for all $i$, and with edge set $\{\{v_{i-1},v_i\}:1\le i\le \ell\}$.
A \emph{partition of $E$ into edge-disjoint paths} means a collection of paths $P_1,\dots,P_t$ in $G$ such that the sets $E(P_i)$ are pairwise disjoint and
\[
E\;=\;\bigsqcup_{i=1}^t E(P_i).
\]
(Vertices may repeat across different paths; only edges are disjoint.)

\smallskip
\noindent\textbf{Conjecture (Erd\H{o}s--Gallai / Gallai path decomposition).}
For every connected graph $G$ on $n$ vertices there exists such a partition with
\[
 t\le \left\lceil \frac n2\right\rceil.
\]
Edge cases: for $n=1$ the bound says $t\le 1$ and $E=\varnothing$ admits $t=0$.
